\section{Кнопочный переключатель}
\begin{verbatim}
					Включение: 1
					Направление: Л 
					Выключение: 2 
					Направление: ПЛ
\end{verbatim}
Код прошивки:
\inputminted{cpp}{z1/main.cpp}
В "a"\ и "b"\ текущее состояние кнопок (инвертированное). При нажатии Левой кнопки - включается, при нажатии Правой а потом Левой - выключается.
\section{RGB светодиод}
\begin{verbatim}
					Фиксация: нет
					Срабатывание: по нажатию 
					Смешивание цветов: да
					Светодиод с общим: анодом
\end{verbatim}
Код прошивки:
\inputminted{cpp}{z2/main.cpp}
В "loop"\ переменная "i"\ приниммает значения от 0 до 2 что позволяет пройтись по всем кнопкам и через вызов "proc"\ зажечь светодиод если кнопка была нажата.
\section{Перетягивание каната}
\begin{verbatim}
					Подключение кнопок: Стягивающий резистор
					Выигрыш: Прерывание по отжатию
					Нажатий для перехода на следующий светодиод: 1
\end{verbatim}
Код прошивки:
\inputminted{cpp}{z3/main.cpp}
В функциях "pushP1"\ и "pushP2"\ сдвигается счёт в лево или право. В основном цикле программы реализовано отображение счёта на светодиодах (счёт инициализирован с двумя светодиодами по центру, чтобы ни у кого небыло преимущества).
\section{Логический симулятор}
\begin{verbatim}
					Логическая функция: 3-input Majority
					Формула: A,B,C (3 кнопки)
					Пояснение: Горит, если нажаты 2 или 3 кнопки.
\end{verbatim}
Код прошивки:
\inputminted{cpp}{z4/main.cpp}
Вводится "a"\ , "b"\  и "c"\ инвертированный сигнал с кнопок, подключенных через встроенный подтягивающий резистор. В запись сигнала на светодиод записано логическое условие через логические операторы.